% =====================================================================
\chapter{文献地图与定位}\label{ch:literature}
% =====================================================================

本章把本文的主张放入文献背景，指出其与既有理论的关系，并标出空缺所在。

\section{六组文献}\label{sec:six-groups}

\paragraph{第一组：现代规格生成与联合评测。}
这是最接近"意图 $\to$ 规格"这一步的前沿。相关工作报告了可测量的结果：对操作系统
内核验证任务的规格生成基准在长上下文与规格完备性上仍显示模型局限
\cite{osvbench2026}；把代码正确性、规格质量与证明生成拆开评估的基准报告了三个维度
之间的显著差距（在特定任务集上，代码正确率、规格健全性与完备性、单次证明成功率
分别为 72.6\%、52.3\% 与 4.9\%）\cite{verina2026}；以可信测试实例化规格的评测
框架指出证明难度会干扰对规格质量的评价\cite{coins2026}；可追溯精化方法把自然语言
要求分解为原子需求并建立需求级追溯\cite{verispecgen2026}；变异驱动方法用语义不等价
变体检验规格强度\cite{specsyn2026}；从代码生成规格的路线应与从意图生成规格的路线
分开研究\cite{specgen2025}。空洞性检测\cite{kupferman2003}则提供了识别"验证通过但
规格空洞成立"的理论路径。这些结果共同支持一个重要判断：\emph{代码生成能力、规格
表达能力与形式化证明能力是不同的研究变量，应当分别评估}。

\paragraph{第二组：程序语义、细化与可观察行为。}
操作语义、指称语义、抽象解释的格模型\cite{cousot1977}、单子语义
\cite{moggi1991,plotkin2002}、域理论、精化关系\cite{hoare1969}与分离逻辑
\cite{appel2007}，共同提供了"何时两个不同抽象层级的系统实现了同一个行为规格"的
数学工具。
语义细化不是简单的信息压缩，而是要定义抽象层之间的\emph{行为关系}。

\paragraph{第三组：类型论、构造性数学与证明提取。}
构造性类型论\cite{hott2013}、证明与程序的对应\cite{howard1980}、以及把证明对象与
构造性计算联系起来的机制，为"AI 提出候选证明、机械内核检查证明"提供了数学基础。
但必须避免一个推论：\emph{如果一个规格已经被翻译成类型，那么规格就必然忠实地表达
了人的意图}。对应关系只保证在选定逻辑与类型系统中命题、类型、证明与项之间的规定
结构对应；人类意图如何转化成命题是另一个需要验证的过程。精化类型
\cite{rondon2008}与 SMT 求解\cite{z32008}是这一方向在工程上的代表；自动定理证明
环境与检索增强的证明搜索\cite{leandojo2023,leancopilot2024}是它在 AI 时代的延续。

\paragraph{第四组：从证明到机器码。}
这是最成熟的一组。经验证编译器把编译转换本身变成被证明的对象\cite{leroy2009}；
经验证编译器的后端覆盖多个目标架构\cite{cakeml2019}；操作系统内核的精化证明展示了
大规模低层系统软件的形式化验证能达到的程度\cite{klein2009}；逐次翻译验证对
\emph{单次}编译结果独立检查\cite{necula2000}；翻译认证为\emph{特定制品}建立可验证
的关系证书\cite{krijnen2024}；证明携带代码让代码提供者附带机器可独立检查的证明
\cite{necula1997}。

这一组的关键教训是：\emph{验证编译器正确性与验证一个特定二进制正确性是两种不同的
策略}。前者证明编译器满足普遍性质，后者对某次产物独立验证。二者可以组合，而不必
把全部信任押在单一工具上。

\paragraph{第五组：资源正确性。}
自动摊还资源分析能在特定语言与成本模型下推导程序的符号化资源上界并具有可靠性证明；
最坏执行时间分析针对特定硬件与假设推导时序上界；静态缓存分析说明硬件模型对上界推断
的影响。必须区分三种时间概念：算法复杂度给出随输入规模变化的成本增长界；最坏执行
时间分析给出特定平台与假设下的上界；实际测量提供特定实验条件下的延迟数据。
\emph{三者不能互相替代}。同样，"线性使用"不等于"最优使用"：一个程序可以只消耗每份
输入资源一次，却仍然使用了不好的算法或引入了过多延迟。

\paragraph{第六组：运行时验证、契约组合与层级涌现。}
这一组包含两类内容，性质不同，必须分开处理。

第一类是\emph{运行时保证与契约}：假设--保证推理要求子系统在输入假设满足时给出相应
保证\cite{henzinger1998}；动态系统的契约组合给出了任意互连拓扑下的组合规则与检查
方法\cite{sharf2021,sharf2024,shali2026}；运行时验证的不确定性研究说明观测、事件
传输与轨迹完整性是监控可靠性的组成部分\cite{taleb2023}；时序逻辑的在线监控
\cite{leucker2009}与信息流安全\cite{goguen1982}分别处理轨迹性质与访问限制。

第二类是\emph{解释框架}：复杂系统的近可分解性为层级结构提供了解释\cite{simon1962}；
最小描述长度为"抽象即压缩"提供了统计与信息论背景
\cite{rissanen1978,grunwald2006,grunwald2019}；计算热力学给出计算的物理下界
\cite{landauer1961}。这一组可以为架构、压缩、模型选择与复杂系统提供解释框架，但
\emph{不能自动提供具体程序的正确性证明}。

\section{按数学功能重排的五层}\label{sec:five-layers}

把上述文献按"它解决链条上哪一段问题"重排，得到一个比"AI 写代码、验证器检查"更深的
结构：

\begin{table}[htbp]
\centering
\small
\caption{按数学功能重排的五个层次}\label{tab:five-layers}
\begin{tabular}{L{1.6cm}L{5.4cm}L{7.2cm}}
\toprule
\textbf{层} & \textbf{功能} & \textbf{代表方向} \\
\midrule
A & 目标的可表达性：形式化规格生成、意图追溯、需求分解、反例生成 &
  规格生成基准、可追溯精化、变异驱动强度检验、空洞性检测、需求形式化 \\
B & 目标的可实现性：逻辑一致性、程序综合、约束求解、可实现性判定 &
  反例引导归纳综合、语法引导综合、程序骨架补全、构造性类型论、SMT 求解 \\
C & 跨层语义保持：操作语义、精化关系、证明检查、编译器正确性、二进制验证 &
  经验证编译器、经验证编译器后端、内核精化证明、证明携带代码、逐次翻译验证 \\
D & 资源与环境保证：复杂度、资源上界、运行时监控、假设--保证契约、执行约束 &
  自动摊还资源分析、最坏执行时间分析、运行时验证、契约组合、信息流安全 \\
E & 结构与物理极限：模型压缩、复杂系统层级、热力学与物理实现 &
  最小描述长度、复杂系统近可分解性、计算热力学、分子自组装 \\
\bottomrule
\end{tabular}
\end{table}

最重要的发现是：这些领域\emph{并没有共同使用一个现成的、完全统一的数学对象}。
但它们已经分别解决了这条链上的不同问题。因此研究机会在于定义它们之间的组合关系，
而不是从零发明全部数学工具。

\section{空缺的交叉点}\label{sec:gaps}

\begin{principle}[文献地图的形状]\label{pr:map-shape}
从规格到二进制的每一步都有成熟文献；从证明到实现的自动化综合有活跃前沿；
但从\emph{意图}到\emph{规格}的鸿沟没有形式化解，只有共识结构，而社会认识论、
拜占庭容错与对抗性评审这些文献存在于完全不同的学科，且尚未与形式化验证系统
系统性地耦合。
\end{principle}

据此，最值得占据的位置是：

\begin{principle}[建议的研究位置]\label{pr:position}
把社会认识论式的对抗性共识与证明携带代码式的机械准入整合在同一个框架内：
让"经共识审定的规格"直接成为"携带证明的制品"的输入，并把共识记录、规格版本与
制品身份用密码学锚定在一起。据本文检视，目前很少有工作同时处理这两个文献群。
\end{principle}

这一位置之所以有价值，是因为它同时满足三个条件：它对应一个真实且尚未解决的技术
缺口（意图层无机械解）；它的两个组成部分都已有成熟的理论基础；它的组合规则可以被
明确写出（定义 \ref{def:two-layer} 与定理 \ref{th:assurance}），因而可以被检验。

\section{与既有理论的关系}\label{sec:relations-lit}

\begin{longtable}{L{3.6cm}L{10.8cm}}
\caption{本文与既有理论的关系}\label{tab:relations}\\
\toprule
\textbf{理论} & \textbf{本文如何使用它，以及不如何用它} \\
\midrule
\endfirsthead
\toprule
\textbf{理论} & \textbf{本文如何使用它，以及不如何用它} \\
\midrule
\endhead
\bottomrule
\endlastfoot
复杂系统的近可分解性 &
  \emph{使用}：作为层级结构之所以稳定的解释，并把它改造成可检验的充分统计量条件
  （定义 \ref{def:neardecomp}、命题 \ref{prop:leakpair}）。
  \emph{不使用}：不作为"所有复杂系统必然具有相同层级"的证据；原文献本身并不证明
  这一点，也不给出验证成本的普遍定理。 \\
抽象解释与 Galois 连接 &
  \emph{使用}：作为层间语义保持的形式工具（定义 \ref{def:galois}）。
  \emph{不使用}：不声称所有层间关系都可用 Galois 连接表达；不能表达时退回到
  模拟关系或人工论证。 \\
最小描述长度 &
  \emph{使用}：作为"理解即压缩"的形式背景与下潜停止准则的来源（定义
  \ref{def:mdl}、命题 \ref{prop:stopping}）。
  \emph{不使用}：不声称最短描述即最正确的描述。MDL 是模型选择原则，其结论依赖所选的
  模型族、编码方式与数据，不能单独决定系统应该实现什么目的。 \\
信息隐藏与深模块 &
  \emph{使用}：作为模块边界的规范性依据。
  \emph{不使用}：本文补充了其可检验版本，即接口是否为行为的充分统计量
  （定义 \ref{def:neardecomp}），因为"信息隐藏得好不好"本身不可测量。 \\
组织结构与系统结构的对应 &
  \emph{使用}：作为高层收敛的机制之一，也是发散的原因之一。
  \emph{不使用}：不作为技术必然性；它是约束的来源，不是解的唯一性来源。 \\
程序逻辑与精化 &
  \emph{使用}：作为规格--实现关系的数学接口。
  \emph{不使用}：不声称安全精化自动保证活性、时限与公平性性质。 \\
证明携带代码 &
  \emph{使用}：作为"生成自由、准入严格"信任结构的经典实现与本文第四支柱的先行者。
  \emph{不使用}：不忽略其前提：接收方必须预先规定安全策略，而策略本身的正确性是
  一个外部问题。 \\
假设--保证推理 &
  \emph{使用}：作为局部保证组合成全局保证的规则来源（命题 \ref{prop:soundness}）。
  \emph{不使用}：不忽略其四个组合条件；组件"单独正确"与系统"整体正确"不互换。 \\
资源分析 &
  \emph{使用}：作为资源正确性可进入形式化分析的证据。
  \emph{不使用}：不把符号上界等同于特定处理器上的时间上界；不把复杂度阶、最坏执行
  时间、统计分位数与实际测量混为一谈。 \\
计算热力学 &
  \emph{使用}：作为"计算的物理代价不可压缩"这一边界条件的依据。
  \emph{不使用}：不用它预测某段普通代码的耗电；不把它并入语义正确性。
  （本文只在第 \ref{ch:objections} 章提及，不作为承重结构。） \\
停机问题与 Rice 定理 &
  \emph{使用}：作为通用综合与通用判定算法不可能存在的边界。
  \emph{不使用}：不用它们证明"意图永远无法形式化"（命题 \ref{prop:no-godel}）。 \\
$\mathrm{P}$ 与 $\mathrm{NP}$ &
  \emph{使用}：只使用 $\mathrm{NP}$ 类的\emph{定义}所蕴含的结构性不对称
  （命题 \ref{prop:asym-struct}）。
  \emph{不使用}：不用"$\mathrm{P} \neq \mathrm{NP}$"或"$\mathrm{P} =
  \mathrm{NP}$"的任何未决断言作为承重结构。 \\
\bottomrule
\end{longtable}

\begin{principle}[本文的定位]\label{pr:self-position}
本文的贡献位于\emph{接口层}，而非新的语义理论：把层级从理解辅助
升级为验证义务的分解单元；把无视集、近可分解性与读集形式化为可计算或可测量的对象；
把内部基础设施化的必然性归因于操作语义的完备性而非代码量；把信任结构写成两种生成
模式的复合并给出共识可靠性的界。这些接口把既有理论连接起来，而不是取代它们。
\end{principle}
